\section*{Hoofdstuk \ref{ch:g2:water-on-earth} --- Water op Aarde}

\begin{solution}{exo:g2:water-on-earth:1}
\omterm{def:g2:water-on-earth:salt-fresh}{Zout water}: de \omterm{def:g2:water-on-earth:ocean}{zee}, de \omterm{def:g2:water-on-earth:ocean}{oceaan}. \omterm{def:g2:water-on-earth:salt-fresh}{Zoet water}: de \omterm{def:g1:landscapes-around-us:water}{rivier}, de regen, het \omterm{def:g1:landscapes-around-us:water}{meer} in
de \omterm{def:g1:landscapes-around-us:landform}{bergen}, het ijs op de \omterm{def:g1:landscapes-around-us:landform}{berg}.
\end{solution}

\begin{solution}{exo:g2:water-on-earth:2}
Omdat het grootste deel van haar oppervlak bedekt is met het blauwe water
van de \omterm{def:g2:water-on-earth:ocean}{oceanen}: vanuit de ruimte gezien ziet ze er blauw uit.
\end{solution}

\begin{solution}{exo:g2:water-on-earth:3}
$100 - 71 = 29$ vierkantjes zijn land.
\end{solution}

\begin{solution}{exo:g2:water-on-earth:4}
Een \omterm{def:g2:water-on-earth:ocean}{oceaan} is groter: een \omterm{def:g2:water-on-earth:ocean}{zee} is een kleinere watervlakte met \omterm{def:g2:water-on-earth:salt-fresh}{zout water},
vaak aan de rand van een \omterm{def:g2:water-on-earth:ocean}{oceaan}.
\end{solution}

\begin{solution}{exo:g2:water-on-earth:5}
Nee: het water van de \omterm{def:g2:water-on-earth:ocean}{oceaan} is \omterm{def:g2:water-on-earth:salt-fresh}{zout water}, en wie het drinkt, verliest
water. De zeeman moet \omterm{def:g2:water-on-earth:salt-fresh}{zoet water} meenemen.
\end{solution}

\begin{solution}{exo:g2:water-on-earth:6}
Bijvoorbeeld: in \omterm{def:g1:landscapes-around-us:water}{rivieren}, in \omterm{def:g1:landscapes-around-us:water}{meren}, in het ijs en de sneeuw van de \omterm{def:g1:landscapes-around-us:landform}{bergen}
en van Antarctica, onder de grond, in de regen.
\end{solution}

\begin{solution}{exo:g2:water-on-earth:7}
\omterm{def:g2:water-on-earth:salt-fresh}{Zout water}: 96 emmers. \omterm{def:g2:water-on-earth:salt-fresh}{Zoet water}: $2 + 2 = 4$ emmers.
(Strikt genomen is een deel van het water onder de grond ook zout; het
plaatje telt het mee bij het zoete water.)
\end{solution}

\begin{solution}{exo:g2:water-on-earth:8}
De \omterm{def:g2:water-on-earth:ocean}{oceanen} (96 emmers), dan het ijs en de sneeuw (2 emmers), dan de
\omterm{def:g1:landscapes-around-us:water}{rivieren} (veel minder dan één emmer).
\end{solution}

\begin{solution}{exo:g2:water-on-earth:9}
Het kan in de grond zakken; het kan over de grond naar een beek of een
\omterm{def:g1:landscapes-around-us:water}{rivier} stromen; het kan weer verdampen als de zon gaat schijnen.
\end{solution}

\begin{solution}{exo:g2:water-on-earth:10}
Op de bodem van de pan blijft zout achter. Het water is verdampt: het is
als waterdamp de lucht in gegaan.
\end{solution}

\begin{solution}{exo:g2:water-on-earth:11}
Bijvoorbeeld: de druppel loopt van het dak in de dakgoot en de afvoer, dan
in een beek en een \omterm{def:g1:landscapes-around-us:water}{rivier}; de \omterm{def:g1:landscapes-around-us:water}{rivier} neemt hem mee naar \omterm{def:g2:water-on-earth:ocean}{zee}; op een zonnige
dag verdampt hij uit de \omterm{def:g2:water-on-earth:ocean}{zee}, stijgt op in de lucht en wordt deel van een
wolk.
\end{solution}
